Adapting large language models (LLMs) to new languages is an expensive and opaque process. Understanding how language models acquire new languages and multilingual abilities is key to achieve efficient adaptation. Prior work on multilingual interpretability research focuses primarily on how trained models process multilingual instructions, leaving unexplored the mechanisms through which they acquire new languages during training. We investigate these training dynamics on decoder-only transformers through the lens of two functional cognitive specializations: language perception (input comprehension) and production (output generation). Through experiments on low-resource languages, we demonstrate how perceptual and productive specialization emerges across different regions of a language model by conducting layer-ablation sweeps along the model's input and output directions. Based on the observed specialization patterns, we propose CogSym, a layer-wise heuristic that enables effective adaptation by exclusively fine-tuning a few early and late layers. We show that tuning only the 25\% outermost layers achieves downstream task performance within 2-3\% deviation from the full fine-tuning baseline. Furthermore, we include an analysis of language learning difficulty in LLMs from the perspective of First Language Acquisition (FLA) and Second Language Acquisition (SLA). Our study provides insights into the resemblance between current LLM architectures and human cognition and behaviour during language acquisition.

\keywords{Large Language Models, Multilingualism, Language Acquisition}